\section*{Lab 3 (k-Means) Numerical Warm-up}
\textbf{Lab 3 focus:} clustering, alternating optimization, initialization sensitivity, scaling, and cluster analysis with k-means.

\medskip
\textbf{Practical usage now:} exploratory grouping, prototype discovery, visual organization, palette compression, and fast structure analysis before downstream labeling.

\medskip
Consider the six unlabeled points in $\mathbb{R}^2$:
\[
  x_1=(1,1),\;
  x_2=(1,2),\;
  x_3=(2,1),\;
  x_4=(6,5),\;
  x_5=(7,5),\;
  x_6=(6,6).
\]
Use $k=2$ with initial centroids
\[
  \mu_1^{(0)}=(1,1),\qquad \mu_2^{(0)}=(7,5).
\]

\begin{enumerate}[label=\textbf{W\arabic*.}]
  \item Perform the assignment step for all six points using squared Euclidean distance.
  \item Compute the updated centroids $\mu_1^{(1)}$ and $\mu_2^{(1)}$.
  \item Compute the objective value
  \[
    J(c,\mu)=\sum_{r=1}^{k}\sum_{i:c_i=r}\|x_i-\mu_r\|_2^2
  \]
  after the update.
  \item Explain why the updated centroids are the arithmetic means of the assigned points.
  \item Suppose the first coordinate were multiplied by $100$ before clustering. Explain how the assignment step would change and why standardization matters.
\end{enumerate}

\section*{Short Conceptual Questions}
\begin{enumerate}[label=\textbf{C\arabic*.}]
  \item State the k-means optimization problem and define every symbol.
  \item Explain why k-means is an alternating minimization algorithm rather than a closed-form one-shot solver.
  \item Explain why different random initializations can lead to different final clusterings.
  \item Give one setting where inertia is useful and one setting where inertia is misleading.
  \item Explain why a PCA plot can help interpret clustering even when clustering is performed in the original feature space.
\end{enumerate}

\section*{Lab 3 (k-Means) Implementation Tasks}
\textbf{Implementation policy (strict):}
\begin{itemize}
  \item Implement the k-means core in NumPy only: assignment, centroid update, inertia computation, and repeated restarts.
  \item You may use plotting libraries for visualization and external datasets/loaders for data access.
  \item Do not call high-level clustering implementations such as \texttt{sklearn.cluster.KMeans} in the core implementation section.
  \item Use the provided starter scaffold in \texttt{ml/assignments/assignment3\_kmeans/}.
  \item Students should primarily edit \texttt{advanced\_ai/clustering/kmeans.py} and use the provided self-checks and task runners.
\end{itemize}

\medskip
\textbf{Core k-means equations to use and reference:}
\begin{itemize}
  \item Within-cluster sum of squares:
  \[
    J(c,\mu)=\sum_{r=1}^{k}\sum_{i:c_i=r}\|x_i-\mu_r\|_2^2.
  \]
  \item Assignment step:
  \[
    c_i=\arg\min_{r \in \{1,\dots,k\}} \|x_i-\mu_r\|_2^2.
  \]
  \item Centroid update:
  \[
    \mu_r=\frac{1}{|\{i:c_i=r\}|}\sum_{i:c_i=r}x_i.
  \]
  \item Standardization using training-set statistics only:
  \[
    x'_j=\frac{x_j-\mu_j}{\sigma_j}.
  \]
\end{itemize}

\medskip
\textbf{Starter architecture to use (required):}
\begin{itemize}
  \item Main worksheet: \texttt{assignment3\_kmeans/kmeans.ipynb}
  \item Internal package: \texttt{assignment3\_kmeans/advanced\_ai/}
  \item Main student edit file: \texttt{advanced\_ai/clustering/kmeans.py}
  \item Self-check script: \texttt{python -m advanced\_ai.self\_checks}
  \item Students should work through the notebook in order and use the code cells as the main implementation guide.
  \item Required task runners: \texttt{iris\_clustering.py}, \texttt{initialization\_scaling\_study.py}, and \texttt{pca\_projection\_clustering.py}
  \item Optional task runner: \texttt{image\_quantization\_optional.py}
  \item Submission notebook: \texttt{assignment3\_kmeans/collect\_submission.ipynb}
  \item Instructor-provided prepared files go in \texttt{assignment3\_kmeans/datasets/}
\end{itemize}

\begin{enumerate}[label=\textbf{K\arabic*.}]
  \item \textbf{Part K1: Build a Reusable k-Means Core}
  \begin{itemize}
    \item \textbf{Goal:} Implement a small NumPy k-means toolkit that supports repeated clustering experiments.
    \item \textbf{What to fill in:}
    \begin{itemize}
      \item random centroid initialization;
      \item nearest-centroid assignment;
      \item centroid update with empty-cluster handling;
      \item inertia/objective tracking;
      \item repeated restarts and best-run selection.
    \end{itemize}
    \item \textbf{Hint:} Verify that inertia never increases within one run.
    \item \textbf{Deliverable:} passing self-checks, toy-data sanity checks, and short notes on how you handled initialization and convergence.
  \end{itemize}

  \item \textbf{Part K2: Clustering the Iris Feature Space}
  \begin{itemize}
    \item \textbf{Public dataset:} Iris dataset (UCI / scikit-learn mirror)\\
    \url{https://archive.ics.uci.edu/dataset/53/iris}\\
    \url{https://scikit-learn.org/stable/auto_examples/datasets/plot_iris_dataset.html}
    \item \textbf{TODO:} Cluster the Iris features using your NumPy k-means implementation.
    \item \textbf{What to fill in:}
    \begin{itemize}
      \item evaluate at least $k \in \{2,3,4,5,6\}$;
      \item run several random restarts for each $k$;
      \item compare raw features and standardized features;
      \item analyze cluster composition against species labels only after clustering is complete.
    \end{itemize}
    \item \textbf{Hint:} Use inertia curves and cluster-composition tables together rather than trusting one number alone.
    \item \textbf{Deliverable:} inertia-vs-$k$ plot, $k=3$ cluster-size summary, one cluster-composition analysis, and a short discussion of how scaling changed the result.
  \end{itemize}

  \item \textbf{Part K3: PCA Projection and Cluster Visualization}
  \begin{itemize}
    \item \textbf{TODO:} Use a provided dataset bundle containing both original features and a 2D PCA projection.
    \item \textbf{What to fill in:}
    \begin{itemize}
      \item cluster the original features with k-means;
      \item visualize the resulting assignments in the 2D PCA plane;
      \item compare clustering behavior in raw space vs standardized space;
      \item explain at least one mismatch between the PCA plot and the final assignments.
    \end{itemize}
    \item \textbf{Hint:} The PCA plot is for interpretation; the clustering objective is still computed in the original feature space you choose.
    \item \textbf{Deliverable:} labeled 2D scatter figure and short interpretation of the observed cluster structure.
  \end{itemize}

  \item \textbf{Part K4: Initialization and Stability Study}
  \begin{itemize}
    \item \textbf{TODO:} Measure how much the final k-means result changes across random initializations.
    \item \textbf{What to fill in:}
    \begin{itemize}
      \item run at least 10 different seeds for one fixed value of $k$;
      \item report the best, median, and worst inertia;
      \item compare the resulting cluster sizes;
      \item explain why local optima appear.
    \end{itemize}
    \item \textbf{Hint:} Keep the dataset fixed and vary only the random initialization.
    \item \textbf{Deliverable:} restart summary table with inertia and cluster sizes, and one short conclusion about why multiple restarts are necessary in practice.
  \end{itemize}

  \item \textbf{Part K5: Optional Advanced Extension}
  \begin{itemize}
    \item Complete \textbf{at most one} optional extension after finishing K1--K4.
  \end{itemize}

  \item \textbf{Part K5A: Image Color Quantization (Optional)}
  \begin{itemize}
    \item \textbf{TODO:} Treat RGB pixel vectors as points in $\mathbb{R}^3$ and cluster them with k-means.
    \item \textbf{Hint:} Replace each pixel by its centroid color and compare the visual result for several values of $k$.
    \item \textbf{Deliverable:} original image, compressed reconstructions for several $k$, and one short note on the trade-off between compression and fidelity.
  \end{itemize}
\end{enumerate}

\section*{Submission Format (Lab 3)}
\begin{itemize}
  \item One PDF report with equations, plots, tables, and conclusions for the warm-up, conceptual questions, and K1--K4.
  \item One completed notebook: \texttt{assignment3\_kmeans/kmeans.ipynb}.
  \item One assignment zip produced from \texttt{assignment3\_kmeans/collect\_submission.ipynb} or \texttt{bash collectSubmission.sh}.
  \item Students may additionally submit \textbf{one optional extension}, but it is not required.
  \item Include random seed, environment file, and reproducibility notes.
  \item For each implementation task: include one short note on what the result says about the k-means objective or its limitations.
  \item Clearly separate exploratory analysis from final conclusions about cluster quality.
\end{itemize}
